\section{Gobierno, pruebas, formación y servicio}
La estimación fija de estas cuentas suma 7.504 HH, de las cuales 2.528 corresponden a implementación y 4.976 a operación. La Tabla~\ref{tab:sd7-gob-serv-curva} agrega sus paquetes. Las sesiones/cohortes, olas y presencia, altas de personal, cobertura externa y correctivos/evolutivos adicionales conservan fórmulas y dimensionamiento propio; no se interpretan como cero por estar fuera de la curva fija.
\begingroup\small\setlength{\tabcolsep}{2pt}
\begin{longtable}{R{9mm}R{12mm}R{12mm}R{12mm}R{12mm}R{12mm}R{12mm}R{14mm}R{12mm}R{14mm}}
\caption{Cargas de gobierno, pruebas, formación y servicio}\label{tab:sd7-gob-serv-curva}\\\hline
 \EncabezadoTabla{Mes} & \EncabezadoTabla{A} & \EncabezadoTabla{D} & \EncabezadoTabla{S} & \EncabezadoTabla{Q} & \EncabezadoTabla{V} & \EncabezadoTabla{I} & \EncabezadoTabla{O+L} & \EncabezadoTabla{J} & \EncabezadoTabla{HH}\\\hline\endfirsthead
\hline \EncabezadoTabla{Mes} & \EncabezadoTabla{A} & \EncabezadoTabla{D} & \EncabezadoTabla{S} & \EncabezadoTabla{Q} & \EncabezadoTabla{V} & \EncabezadoTabla{I} & \EncabezadoTabla{O+L} & \EncabezadoTabla{J} & \EncabezadoTabla{HH}\\\hline\endhead
1 & 16 & 0 & 0 & 8 & 0 & 8 & 0 & 32 & 64 \\\hline
2 & 24 & 0 & 16 & 32 & 0 & 8 & 0 & 56 & 136\\\hline
3 & 8 & 8 & 8 & 8 & 0 & 16 & 0 & 32 & 80 \\\hline
4 & 8 & 0 & 0 & 16 & 0 & 8 & 0 & 40 & 72\\\hline
5 & 8 & 0 & 0 & 8 & 0 & 0 & 0 & 16 & 32 \\\hline
6 & 8 & 0 & 8 & 8 & 0 & 8 & 0 & 24 & 56\\\hline
7 & 8 & 0 & 0 & 8 & 0 & 0 & 0 & 16 & 32 \\\hline
8 & 8 & 0 & 0 & 8 & 0 & 0 & 0 & 16 & 32\\\hline
9 & 16 & 0 & 0 & 24 & 0 & 8 & 0 & 16 & 64 \\\hline
10 & 8 & 0 & 0 & 32 & 8 & 8 & 0 & 16 & 72\\\hline
11 & 80 & 24 & 40 & 184 & 16 & 48 & 160 & 24 & 576 \\\hline
12 & 24 & 0 & 4 & 32 & 0 & 8 & 0 & 28 & 96\\\hline
13 & 8 & 0 & 0 & 8 & 0 & 0 & 0 & 16 & 32 \\\hline
14 & 8 & 0 & 0 & 8 & 0 & 0 & 0 & 16 & 32\\\hline
15 & 32 & 8 & 8 & 48 & 0 & 16 & 24 & 32 & 168 \\\hline
16 & 8 & 0 & 0 & 8 & 0 & 0 & 0 & 16 & 32\\\hline
17 & 80 & 24 & 40 & 208 & 24 & 56 & 160 & 24 & 616 \\\hline
18 & 24 & 0 & 12 & 32 & 0 & 16 & 16 & 44 & 144\\\hline
19 & 8 & 0 & 0 & 8 & 0 & 0 & 0 & 16 & 32 \\\hline
20 & 24 & 8 & 8 & 40 & 0 & 8 & 48 & 24 & 160\\\hline
21 & 0 & 8 & 12 & 24 & 8 & 12 & 44 & 4 & 112 \\\hline
22 & 0 & 8 & 12 & 24 & 8 & 12 & 44 & 4 & 112\\\hline
23 & 0 & 8 & 12 & 24 & 8 & 12 & 44 & 4 & 112 \\\hline
24 & 8 & 16 & 24 & 32 & 8 & 28 & 60 & 16 & 192\\\hline
25 & 0 & 8 & 12 & 24 & 8 & 12 & 44 & 4 & 112 \\\hline
26 & 0 & 24 & 20 & 48 & 8 & 36 & 84 & 4 & 224\\\hline
27 & 0 & 8 & 12 & 20 & 8 & 8 & 28 & 4 & 88 \\\hline
28 & 0 & 8 & 12 & 20 & 8 & 8 & 28 & 4 & 88\\\hline
29 & 0 & 8 & 12 & 20 & 8 & 8 & 28 & 4 & 88 \\\hline
30 & 8 & 8 & 16 & 28 & 8 & 16 & 44 & 8 & 136\\\hline
31 & 0 & 8 & 12 & 20 & 8 & 8 & 28 & 4 & 88 \\\hline
32 & 0 & 24 & 52 & 60 & 8 & 40 & 68 & 4 & 256\\\hline
33 & 0 & 8 & 12 & 20 & 8 & 8 & 28 & 4 & 88 \\\hline
34 & 0 & 8 & 12 & 20 & 8 & 8 & 28 & 4 & 88\\\hline
35 & 0 & 8 & 12 & 20 & 8 & 8 & 28 & 4 & 88 \\\hline
36 & 8 & 16 & 24 & 28 & 8 & 24 & 44 & 16 & 168\\\hline
37 & 0 & 8 & 12 & 20 & 8 & 8 & 28 & 4 & 88 \\\hline
38 & 0 & 24 & 20 & 44 & 8 & 32 & 68 & 4 & 200\\\hline
39 & 0 & 8 & 12 & 20 & 8 & 8 & 28 & 4 & 88 \\\hline
40 & 0 & 8 & 12 & 20 & 8 & 8 & 28 & 4 & 88\\\hline
41 & 0 & 8 & 12 & 20 & 8 & 8 & 28 & 4 & 88 \\\hline
42 & 8 & 8 & 16 & 28 & 8 & 16 & 44 & 8 & 136\\\hline
43 & 0 & 8 & 12 & 20 & 8 & 8 & 28 & 4 & 88 \\\hline
44 & 0 & 24 & 52 & 60 & 8 & 40 & 68 & 4 & 256\\\hline
45 & 0 & 8 & 12 & 20 & 8 & 8 & 28 & 4 & 88 \\\hline
46 & 0 & 8 & 12 & 20 & 8 & 8 & 28 & 4 & 88\\\hline
47 & 0 & 8 & 12 & 20 & 8 & 8 & 28 & 4 & 88 \\\hline
48 & 8 & 16 & 24 & 28 & 8 & 24 & 44 & 16 & 168\\\hline
49 & 0 & 8 & 12 & 20 & 8 & 8 & 28 & 4 & 88 \\\hline
50 & 0 & 24 & 20 & 44 & 8 & 32 & 68 & 4 & 200\\\hline
51 & 0 & 8 & 12 & 20 & 8 & 8 & 28 & 4 & 88 \\\hline
52 & 0 & 8 & 12 & 20 & 8 & 8 & 28 & 4 & 88\\\hline
53 & 0 & 56 & 28 & 36 & 8 & 48 & 44 & 12 & 232 \\\hline
54 & 8 & 32 & 48 & 60 & 16 & 64 & 92 & 16 & 336\\\hline
55 & 0 & 16 & 12 & 28 & 8 & 24 & 44 & 4 & 136 \\\hline
56 & 0 & 32 & 52 & 68 & 8 & 56 & 84 & 4 & 304\\\hline
\end{longtable}\FuenteTabla{Paquetes fijos; excluye sesiones por cohorte, presencia por ola, cobertura externa, altas de personal y trabajo correctivo variable.}\endgroup


La suma parcial fija de implementación alcanza 13.588 HH al agregar estas cuentas a las 11.060 anteriores; operación fija alcanza 7.704 HH con estabilización de datos, piloto estacional IN2 y seguimiento fijo de innovación. Se agregan revisión documental, enrolamiento, sesiones, presencia por ola, 448 HH de estabilización presencial E1 y 448 de E2, seguimiento IN2, revisiones de botaduras y cobertura externa. El componente estacional requiere meses reales, aunque ya tenga esfuerzo fijo. Los informes técnicos no se convierten en cobertura NOC/SOC o acompañamiento en terreno.

Las ventanas propuestas para pruebas y recepción mantienen H4/H5 y H9/H10, pero deben integrar liberaciones, observaciones, correcciones, reensayos y disponibilidad de QA/Seguridad/SRE. La prueba UAT final depende de las demás pruebas y de los ensayos de migración sin incluirse como su propia predecesora. La secuencia completa no está demostrada por agrupar todos los ensayos en un mes; PERT/CPM deberá cuantificarla y contrastarla con el calendario protegido. Las comprobaciones recurrentes se anticiparán si una fecha relativa cae en congelamiento, conservando intervalo máximo semestral/anual y aprobación de la ventana.

La salida conserva acompañamiento real de al menos 90 días dentro de 53--56 y repositorio/credenciales/exportaciones recibidos antes de eliminar un origen. Se mantiene el plan inicial dentro de 90 días y su actualización anual. Custodia independiente de fuentes, renovación de certificaciones y servicios/licencias de todo el período se registrarán en sus contratos y habilitaciones, sin afirmar que un informe equivale a custodia contratada.

La nivelación completa sumará por persona todos los paquetes fijos y variables, cobertura/otros compromisos y ausencias. También instanciará las funciones de refuerzo y proveedores externos, con ubicación, competencias, turnos, disponibilidad y costo. La capacidad general de servicio no se considera demostrada mientras esos datos y el volumen de trabajo no estén incorporados; la etapa de diseño permite fijar el modelo, pero no atribuir contratación o verificación ejecutada.

\begin{longtable}{R{18mm}L{36mm}R{30mm}R{30mm}}
\caption{Balance conjunto de paquetes fijos con mes asignado}\label{tab:sd7-balance-final-fijo}\\\hline
 \EncabezadoTabla{Mes} & \EncabezadoTabla{Perfil/persona} & \EncabezadoTabla{HH conocidas} & \EncabezadoTabla{Exceso sobre 144}\\\hline\endfirsthead
\hline \EncabezadoTabla{Mes} & \EncabezadoTabla{Perfil/persona} & \EncabezadoTabla{HH conocidas} & \EncabezadoTabla{Exceso sobre 144}\\\hline\endhead
3 & A & 352 & 208 \\\hline
3 & D & 256 & 112\\\hline
6 & O+L & 160 & 16 \\\hline
8 & D & 304 & 160\\\hline
8 & Q & 192 & 48 \\\hline
8 & V & 1208 & 1064\\\hline
9 & Q & 160 & 16 \\\hline
9 & V & 840 & 696\\\hline
11 & Q & 244 & 100 \\\hline
11 & O+L & 176 & 32\\\hline
13 & A & 232 & 88 \\\hline
13 & D & 152 & 8\\\hline
15 & V & 640 & 496 \\\hline
16 & Q & 188 & 44\\\hline
16 & V & 656 & 512 \\\hline
17 & Q & 364 & 220\\\hline
17 & V & 408 & 264 \\\hline
17 & O+L & 176 & 32\\\hline
18 & Q & 200 & 56 \\\hline
20 & Q & 192 & 48\\\hline
21 & O+L & 192 & 48 \\\hline
\end{longtable}\FuenteTabla{Todos los paquetes fijos con mes asignado; excluye componentes estacionales sin mes efectivo, variables, presencia y cobertura externa. No constituye nivelación aprobada.}

El balance conjunto sustituye, para el alcance fijo con mes asignado, los contrastes parciales anteriores. Las sobrecargas permanecen visibles y requieren solución; no se ofrecen los mismos titulares simultáneamente a tareas incompatibles. Aun después de corregirlas se deberá comprobar la quincena, la duración efectiva, el calendario protegido y los hitos de recepción.
